\section{Conclusiones y Próximos Pasos}
\label{sec:conclusiones}

El incremento de este ciclo culmina con la solidificación de la arquitectura técnica y una integración de interfaces altamente optimizada, permitiendo que la interacción del usuario final fluya sobre una capa de datos robusta, estricta y segura.

\subsection{Revisión de la Definition of Done (DoD)}
\begin{ChecklistDoD}
    \item Esquema lógico desplegado en Neon.tech con nomenclatura en inglés y tipos restrictivos.
    \item Integración del flujo de autenticación y gestión de sesiones mediante OAuth2 certificada.
    \item Vistas responsivas, flujos de validación cronológica y herramientas de conversión implementadas en UI.
\end{ChecklistDoD}

\vspace{0.5cm}

\begin{AlertaDefecto}
\textbf{Integración Pendiente en Registro Interno:} Las APIs transaccionales del registro de usuarios y los componentes visuales del formulario se desarrollaron de manera exitosa e independiente en este sprint. Su conexión y enrutamiento \textit{End-to-End} se consolida como la primera prioridad para el inicio del próximo ciclo.
\end{AlertaDefecto}

El aplazamiento planificado en la selección del tipo de usuario (dirigiéndolo hacia la edición del perfil) demostró ser una decisión arquitectónica acertada para minimizar la fricción inicial y aumentar la tasa de retención. La base técnica establecida reduce drásticamente la deuda técnica para las historias de usuario subsecuentes.